\section{Discussion}
\label{sec:discussion}

\subsection{Soundness vs.\ Completeness Tradeoffs}

The results expose a fundamental tension in the design of equational implication
classifiers.  The BFS reachability search (step~(3) of the classifier)
is fast and returns high-confidence scores for most implication cases, but its
bounded character means it may occasionally produce false positives: two terms whose
reachable sets collide by coincidence, without a true equational proof.  KB completion
is slower and incomplete (it times out for complex laws), but whenever it terminates
successfully, it provides a machine-checkable proof of implication with no false
positives.

The two-tier architecture---KB first, BFS fallback---reflects this tradeoff explicitly.
Approximately $12\%$ of implication cases receive KB-backed proofs; the remaining $88\%$
rely on BFS.  Increasing the KB timeout (from 0.8~s to 2--3~s with multiprocessing)
would likely raise the KB coverage fraction and reduce the false-positive risk from BFS.

Similarly, Z3 at fixed domain sizes $n \leq 5$ is complete for those sizes but cannot
detect infinite-magma counterexamples.  By Remark~\ref{rem:infinite} and Austin's
theorem \cite{austin1979equational}, a law pair $E, E'$ may satisfy $E \models_\mathrm{fin} E'$
(no finite counterexample) while having $E \not\models E'$ (an infinite counterexample
exists).  For such pairs, \textsc{predictor4} would incorrectly assign high implication
probability.  The ETP ground truth (which includes infinite models) captures these cases;
our evaluation is against this complete ground truth.

\subsection{Comparison with Prior Approaches}

Table~\ref{tab:comparison} summarizes the components of \textsc{predictor3} and
\textsc{predictor4} and their theoretical properties.

\begin{table}[ht]
\centering
\caption{Component comparison between \textsc{predictor3} and \textsc{predictor4}.}
\label{tab:comparison}
\renewcommand{\arraystretch}{1.15}
\begin{tabular}{@{}lp{3.2cm}p{3.2cm}l@{}}
\toprule
\textbf{Component} & \textbf{\textsc{predictor3}} & \textbf{\textsc{predictor4}} & \textbf{Status} \\
\midrule
Positive proof & BFS (bounded, unsound in pathological cases)
               & KB completion + BFS fallback
               & Improved \\
Neg.\ search ($n \leq 3$) & Exhaustive (sound) & Exhaustive (sound) & Same \\
Neg.\ search ($n = 4, 5$) & Hill-climbing (heuristic)
               & Z3 SMT (complete per $n$)
               & Improved \\
Algebraic tests & Polynomial evaluation & Polynomial evaluation & Same \\
Fallback & Structural heuristic & Structural heuristic & Same \\
\bottomrule
\end{tabular}
\end{table}

The ETP itself used Vampire, Prover9, MagmaEgg, and Mace4 for its verification
\cite{etp2025}---industrial-strength tools unavailable as Python libraries.
\textsc{predictor4} implements KB completion and SMT model finding from scratch in
Python, intentionally staying within a self-contained Python environment suitable for
rapid prototyping and extension.

\subsection{The Soundness Invariant in Context}

The validity condition $\operatorname{vars}(r) \subseteq \operatorname{vars}(\ell)$ is standard in term
rewriting textbooks \cite{baader1999term} but is easy to violate in practice when
variable renaming for critical pair computation is implemented carelessly.  A correct
renaming procedure ensures that the variables of the two rules being superposed are
disjoint \emph{before} computing the unifier, and that no variable name from either
rule leaks into the renamed copies.  The bug in the original implementation
(Example~\ref{ex:bug}) was precisely such a naming leak: a renamed variable $y_0$
appeared in the RHS of a critical pair without being bound by the LHS.

Our fix---checking $\operatorname{vars}(r) \subseteq \operatorname{vars}(\ell)$ before accepting any rule---is
a defensive measure that acts as a post-hoc safety net.  A correct implementation
should not need this check; its presence signals a defect in variable renaming.
Nonetheless, the check has value as a runtime assertion that guards against similar
bugs in future modifications.

\subsection{Open Questions}

Several questions arise from this work.

\begin{enumerate}[(1)]
  \item \textbf{KB termination for magma laws.}
        For what classes of single-equation magma systems does KB completion
        terminate?  Associativity is a known divergence case (KB produces an
        infinite TRS of left- or right-association rules).  Are there algebraic
        characterizations of laws for which KB terminates?

  \item \textbf{Proof certificates.}
        A convergent TRS $R$ produced by KB, together with the critical pair
        join sequence, constitutes a machine-checkable proof of implication
        \cite{sternagel2012recording}.  Can \textsc{predictor4} be extended to
        export Lean~4 proof certificates compatible with the ETP formalization?

  \item \textbf{Larger domain sizes for Z3.}
        Extending Z3 model finding to $n = 6, 7$ would increase coverage of
        non-implications with larger minimal counterexamples.  The constraint
        size grows as $n^{n^2 + |\operatorname{vars}|}$; empirical investigation of the
        Z3 timeout boundary as a function of $n$ and law complexity is needed.

  \item \textbf{Equivalence class caching.}
        The 4,694 laws form 1,415 equivalence classes under the implication preorder.
        Laws in the same equivalence class have identical implication relations.
        KB results cached per equivalence class could amortize the KB timeout
        cost significantly.

  \item \textbf{Better reduction orders.}
        LPO with fixed variable precedence is one reduction order; many alternatives
        exist (Knuth-Bendix order, polynomial interpretations).  Trying multiple
        orderings when LPO fails to orient an equation would increase KB convergence
        rates.
\end{enumerate}
